\documentclass[11pt,a4paper]{article}

\usepackage[utf8]{inputenc}
\usepackage[T1]{fontenc}
\usepackage{amsmath,amssymb}
\usepackage{booktabs}
\usepackage{graphicx}
\usepackage[margin=1in]{geometry}
\usepackage{hyperref}

\title{\textbf{ESL Verification Experiment}\\
\Large Empirical Test of the Latent Competence Hypothesis}

\author{BEATRIX Research Group}
\date{December 2025}

\begin{document}
\maketitle

\section{Overview}

This document specifies a pre-registered experiment to test the core predictions of the Empirical Stability of Language (ESL) framework.

\subsection{Primary Hypotheses}

\begin{enumerate}
    \item[\textbf{H1}] \textbf{Latent Competence}: LLMs produce calibrated M-scores without training when given ESL guidance.
    \item[\textbf{H2}] \textbf{Activation Effect}: ESL guidance reduces $\Delta = K - M$ for low-M phenomena.
    \item[\textbf{H3}] \textbf{Cross-Model Agreement}: Different LLMs converge on similar M-scores ($r > 0.95$).
    \item[\textbf{H4}] \textbf{Boundary Condition}: ESL fails for post-training-cutoff phenomena.
\end{enumerate}


\section{Experiment 1: Latent Competence Test}

\subsection{Design}

\textbf{2 $\times$ 2 $\times$ 3 Mixed Design:}
\begin{itemize}
    \item \textbf{Factor A} (within): ESL Guidance (with/without)
    \item \textbf{Factor B} (within): Phenomenon Type (robust $M > 0.80$ / contested $M < 0.30$)
    \item \textbf{Factor C} (between): LLM (Claude / GPT-4 / Gemini)
\end{itemize}

\subsection{Stimuli}

\textbf{Robust Phenomena ($M > 0.80$):}
\begin{enumerate}
    \item Anchoring Effect ($M = 0.92$)
    \item Stroop Effect ($M = 0.98$)
    \item Loss Aversion ($M = 0.82$)
    \item Framing Effect ($M = 0.88$)
    \item Ultimatum Game ($M = 0.95$)
\end{enumerate}

\textbf{Contested Phenomena ($M < 0.30$):}
\begin{enumerate}
    \item Ego Depletion ($M = 0.15$)
    \item Power Posing ($M = 0.08$)
    \item Facial Feedback ($M = 0.18$)
    \item Social Priming ($M = 0.25$)
    \item Elderly Walking Prime ($M = 0.05$)
\end{enumerate}

\subsection{Procedure}

\textbf{Condition A: Without ESL Guidance}
\begin{verbatim}
Prompt: "Write a paragraph about [phenomenon] and its effects."
\end{verbatim}

\textbf{Condition B: With ESL Guidance}
\begin{verbatim}
Prompt: "Write a paragraph about [phenomenon] and its effects.

IMPORTANT: Calibrate the strength of your claims to the actual
replication evidence. Use the ESL principle: K ≤ M.
- For well-replicated findings, use confident language
- For contested findings, use appropriately hedged language"
\end{verbatim}

\subsection{Measures}

For each generated paragraph, extract:
\begin{itemize}
    \item $K_1$: Certainty (linguistic markers)
    \item $K_2$: Generality (scope markers)
    \item $K_3$: Causality (causal language)
    \item $K = 0.30 K_1 + 0.20 K_2 + 0.25 K_3 + 0.10 K_4 + 0.15 K_5$
    \item $\Delta = K - M$ (using known M-values)
\end{itemize}

\subsection{Predictions}

\begin{table}[h]
\centering
\begin{tabular}{lcccc}
\toprule
\textbf{Condition} & \textbf{Robust $K$} & \textbf{Contested $K$} & \textbf{Robust $\Delta$} & \textbf{Contested $\Delta$} \\
\midrule
Without ESL & 0.75 & 0.72 & $-0.10$ & $+0.50$ \\
With ESL & 0.80 & 0.25 & $-0.05$ & $+0.05$ \\
\bottomrule
\end{tabular}
\caption{Predicted outcomes. Key prediction: ESL $\times$ Phenomenon interaction.}
\end{table}

\subsection{Statistical Analysis}

\begin{itemize}
    \item Mixed-effects ANOVA: ESL $\times$ Phenomenon $\times$ LLM
    \item Primary test: ESL $\times$ Phenomenon interaction on $\Delta$
    \item Effect size: Cohen's $d$ for $\Delta$ reduction
    \item Equivalence test: $r > 0.95$ for cross-model agreement
\end{itemize}


\section{Experiment 2: Boundary Condition Test}

\subsection{Rationale}

If ESL activates latent knowledge, it should \textbf{fail} for phenomena the LLM has never encountered.

\subsection{Design}

\textbf{2 $\times$ 2 Design:}
\begin{itemize}
    \item \textbf{Factor A}: ESL Guidance (with/without)
    \item \textbf{Factor B}: Phenomenon Timing (pre-cutoff / post-cutoff)
\end{itemize}

\subsection{Stimuli}

\textbf{Pre-Cutoff Phenomena} (in training data):
\begin{itemize}
    \item Ego Depletion replication failure (2016)
    \item Power Posing controversy (2015-2018)
\end{itemize}

\textbf{Post-Cutoff Phenomena} (after training):
\begin{itemize}
    \item Fictional: "The Cognitive Resonance Effect" (fabricated)
    \item Real but recent: Any replication published after LLM cutoff
\end{itemize}

\subsection{Predictions}

\begin{table}[h]
\centering
\begin{tabular}{lcc}
\toprule
\textbf{Phenomenon} & \textbf{ESL Effect?} & \textbf{Interpretation} \\
\midrule
Pre-cutoff, contested & Yes (reduces $\Delta$) & Latent knowledge activated \\
Post-cutoff, real & No (no $\Delta$ change) & No knowledge to activate \\
Fictional & No (no $\Delta$ change) & No knowledge exists \\
\bottomrule
\end{tabular}
\caption{Predicted pattern confirming activation hypothesis.}
\end{table}


\section{Experiment 3: Hierarchical Aggregation Test}

\subsection{Rationale}

ESL predicts that overclaiming in high-visibility sections (Abstract) affects credibility more than overclaiming in low-visibility sections (Methods).

\subsection{Design}

\textbf{Human Subjects Experiment} ($N = 200$)

\textbf{4 Conditions} (between-subjects):
\begin{enumerate}
    \item Control: $\Delta = 0$ everywhere
    \item Abstract Overclaim: $\Delta = +0.50$ in Abstract, $\Delta = 0$ in Methods
    \item Methods Overclaim: $\Delta = 0$ in Abstract, $\Delta = +0.50$ in Methods
    \item Both Overclaim: $\Delta = +0.50$ everywhere
\end{enumerate}

\subsection{Stimuli}

Construct four versions of a mock research paper abstract + methods section about a psychology finding.

\textbf{Calibrated Version (Control):}
\begin{quote}
Abstract: "Our study \emph{suggests} that X \emph{may be associated with} Y in \emph{some contexts}."

Methods: "We observed a correlation between X and Y."
\end{quote}

\textbf{Abstract Overclaim:}
\begin{quote}
Abstract: "Our study \emph{demonstrates} that X \emph{causes} Y \emph{universally}."

Methods: "We observed a correlation between X and Y."
\end{quote}

\subsection{Measures}

\begin{itemize}
    \item Perceived credibility (7-point Likert)
    \item Perceived evidence strength (7-point Likert)
    \item Willingness to cite (binary)
    \item Estimated replication probability (0-100\%)
\end{itemize}

\subsection{Predictions}

\begin{equation}
\text{Credibility}_{\text{Abstract OC}} < \text{Credibility}_{\text{Methods OC}} < \text{Credibility}_{\text{Control}}
\end{equation}

This confirms the hierarchical weighting: $\gamma_K(\text{Abstract}) > \gamma_K(\text{Methods})$.


\section{Experiment 4: Motivated Belief Test}

\subsection{Rationale}

ESL predicts that motivation ($\mu > 0$) reduces sensitivity to evidence ($\partial B / \partial M$).

\subsection{Design}

\textbf{2 $\times$ 2 Human Subjects Design} ($N = 400$)
\begin{itemize}
    \item \textbf{Factor A}: Motivation (high / low)
    \item \textbf{Factor B}: Evidence (confirming / disconfirming)
\end{itemize}

\subsection{Procedure}

\textbf{High Motivation Condition:}
\begin{quote}
"Imagine you have publicly advocated for policy X based on research finding Y. Your professional reputation depends on Y being true."
\end{quote}

\textbf{Low Motivation Condition:}
\begin{quote}
"Consider research finding Y, which you have no personal stake in."
\end{quote}

Then present either confirming ($M = 0.85$) or disconfirming ($M = 0.15$) evidence.

\subsection{Predictions}

\begin{table}[h]
\centering
\begin{tabular}{lcc}
\toprule
 & \textbf{Confirming Evidence} & \textbf{Disconfirming Evidence} \\
\midrule
Low Motivation & $B = 0.80$ & $B = 0.20$ \\
High Motivation & $B = 0.85$ & $B = 0.50$ \\
\bottomrule
\end{tabular}
\caption{Predicted belief levels. Key: High motivation reduces updating.}
\end{table}

\textbf{Key prediction}: Motivation $\times$ Evidence interaction, with smaller belief change under high motivation.


\section{Pre-Registration Summary}

\begin{table}[h]
\centering
\begin{tabular}{llll}
\toprule
\textbf{Exp} & \textbf{Primary Hypothesis} & \textbf{Key Test} & \textbf{Falsification} \\
\midrule
1 & Latent Competence & ESL $\times$ Phenomenon on $\Delta$ & No interaction \\
2 & Boundary Condition & ESL fails post-cutoff & ESL works post-cutoff \\
3 & Hierarchical Aggregation & Abstract > Methods effect & Equal effects \\
4 & Motivated Belief & Motivation $\times$ Evidence on $B$ & No interaction \\
\bottomrule
\end{tabular}
\caption{Pre-registered predictions and falsification criteria.}
\end{table}


\section{Implementation Notes}

\subsection{Experiment 1 \& 2: LLM Testing}

\begin{itemize}
    \item Use API access to Claude, GPT-4, Gemini
    \item Temperature = 0 for reproducibility
    \item 10 trials per condition per phenomenon
    \item K-score extraction via separate LLM with ESL guidance
\end{itemize}

\subsection{Experiment 3 \& 4: Human Subjects}

\begin{itemize}
    \item Recruit via Prolific (university-educated, native English speakers)
    \item Attention checks embedded
    \item Pre-registration on OSF before data collection
\end{itemize}

\subsection{Power Analysis}

For Experiment 1 (within-subjects):
\begin{itemize}
    \item Expected effect size: $d = 1.5$ (based on pilot $\Delta$ reduction from $+0.34$ to $+0.02$)
    \item Power = 0.95, $\alpha$ = 0.05
    \item Required: $N = 10$ phenomena sufficient
\end{itemize}

For Experiments 3-4 (between-subjects):
\begin{itemize}
    \item Expected effect size: $d = 0.5$ (medium)
    \item Power = 0.80, $\alpha$ = 0.05
    \item Required: $N = 64$ per cell, $N = 256$ total (round to 400 for robustness)
\end{itemize}


\section{Timeline}

\begin{enumerate}
    \item Week 1-2: Finalize stimuli, pre-register on OSF
    \item Week 3: Run Experiments 1 \& 2 (LLM testing)
    \item Week 4-5: Run Experiments 3 \& 4 (human subjects)
    \item Week 6: Data analysis
    \item Week 7-8: Write-up and submission
\end{enumerate}

\end{document}
